% ============================================================
%  CLASE: Sistemas de ecuaciones lineales con dos incógnitas
%  Aplicaciones en Agronomía — Nivel universitario
%  Compilar con: pdflatex sistemas_lineales_agronomia.tex
% ============================================================

\documentclass[11pt, aspectratio=169]{beamer}

% ---------- Tema y colores agronomía ----------
\usetheme{Madrid}

\definecolor{verdecampo}{RGB}{34, 100, 34}
\definecolor{verdeagro}{RGB}{74, 145, 50}
\definecolor{tierracolor}{RGB}{139, 90, 43}
\definecolor{celesteagua}{RGB}{30, 130, 190}
\definecolor{amarillotrigo}{RGB}{210, 170, 30}
\definecolor{grisclaro}{RGB}{245, 246, 242}
\definecolor{rojoalerta}{RGB}{175, 35, 35}

\setbeamercolor{palette primary}{bg=verdecampo, fg=white}
\setbeamercolor{palette secondary}{bg=verdeagro, fg=white}
\setbeamercolor{palette tertiary}{bg=verdecampo!80, fg=white}
\setbeamercolor{palette quaternary}{bg=verdecampo, fg=white}
\setbeamercolor{structure}{fg=verdecampo}
\setbeamercolor{block title}{bg=verdecampo, fg=white}
\setbeamercolor{block body}{bg=grisclaro, fg=black}
\setbeamercolor{block title alerted}{bg=rojoalerta, fg=white}
\setbeamercolor{block body alerted}{bg=rojoalerta!10, fg=black}
\setbeamercolor{block title example}{bg=verdeagro, fg=white}
\setbeamercolor{block body example}{bg=verdeagro!10, fg=black}
\setbeamercolor{frametitle}{bg=verdecampo, fg=white}
\setbeamercolor{title}{bg=verdecampo, fg=white}

% ---------- Fuentes ----------
\usefonttheme{professionalfonts}
\usepackage[T1]{fontenc}
\usepackage[utf8]{inputenc}
\usepackage[spanish]{babel}

% ---------- Paquetes matemáticos ----------
\usepackage{amsmath, amssymb, amsfonts}
\usepackage{mathtools}
\usepackage{array}
\usepackage{booktabs}
\usepackage{multicol}
\usepackage{tikz}
\usetikzlibrary{intersections, arrows.meta, calc}
\usepackage{pgfplots}
\pgfplotsset{compat=1.18}

% ---------- Otros paquetes ----------
\usepackage{xcolor}
\usepackage{tcolorbox}
\tcbuselibrary{theorems, skins, breakable}
\usepackage{siunitx}

% ---------- Navegación ----------
\setbeamertemplate{navigation symbols}{}
\setbeamertemplate{footline}{
  \leavevmode
  \hbox{
    \begin{beamercolorbox}[wd=.35\paperwidth, ht=2.4ex, dp=1ex, center]{palette primary}
      \usebeamerfont{author in head/foot}\insertshortauthor
    \end{beamercolorbox}
    \begin{beamercolorbox}[wd=.45\paperwidth, ht=2.4ex, dp=1ex, center]{palette secondary}
      \usebeamerfont{title in head/foot}\insertshorttitle
    \end{beamercolorbox}
    \begin{beamercolorbox}[wd=.2\paperwidth, ht=2.4ex, dp=1ex, right]{palette primary}
      \usebeamerfont{date in head/foot}
      \insertframenumber{} / \inserttotalframenumber\hspace*{1.5ex}
    \end{beamercolorbox}
  }
}

% ---------- Metadatos ----------
\title[Sistemas Lineales — Agronomía]{Sistemas de Ecuaciones Lineales\\con Dos Incógnitas}
\subtitle{Matemáticas Aplicadas a la Agronomía — Unidad 2}
\author[Dept. Ciencias Básicas]{Departamento de Ciencias Básicas}
\institute{Facultad de Agronomía}
\date{\today}

% ---------- Comandos útiles ----------
\newcommand{\R}{\mathbb{R}}
\newcommand{\mat}[1]{\begin{pmatrix} #1 \end{pmatrix}}
\newcommand{\deter}[4]{\begin{vmatrix} #1 & #2 \\ #3 & #4 \end{vmatrix}}
\newcommand{\unidad}[1]{\,\text{#1}}

% ===========================================================
\begin{document}
% ===========================================================

% ---- Portada -----------------------------------------------
\begin{frame}[plain]
  \titlepage
\end{frame}

% ---- Índice ------------------------------------------------
\begin{frame}{Contenidos}
  \tableofcontents
\end{frame}

% ===========================================================
\section{Motivación agronómica}
% ===========================================================

\begin{frame}{¿Por qué sistemas de ecuaciones en agronomía?}
  \begin{columns}[T]
    \column{0.55\textwidth}
      En la práctica agrícola, es frecuente tener \textbf{dos condiciones simultáneas}
      que determinan una situación óptima:
      \begin{itemize}
        \item Dosis de \textbf{fertilizantes} con restricciones de N y P
        \item \textbf{Mezcla de plaguicidas} con concentración objetivo
        \item Balance \textbf{hídrico} en sistemas de riego
        \item Asignación de \textbf{superficie} entre cultivos
        \item Composición de \textbf{sustratos} y enmiendas
      \end{itemize}
    \column{0.42\textwidth}
      \begin{exampleblock}{Pregunta motivadora}
        Un agricultor dispone de dos fertilizantes:
        \begin{itemize}
          \item Fertilizante A: 20\% N, 10\% P
          \item Fertilizante B: 5\% N, 15\% P
        \end{itemize}
        ¿Cuántos kg de cada uno aplicar para suministrar exactamente
        \textbf{30 kg/ha de N} y \textbf{20 kg/ha de P}?
      \end{exampleblock}
  \end{columns}
\end{frame}

% ===========================================================
\section{Definición y forma general}
% ===========================================================

\begin{frame}{Definición general}
  \begin{block}{Sistema lineal $2\times2$}
    Dados los parámetros agronómicos $(a_i, b_i, c_i \in \R)$,
    un sistema lineal con dos incógnitas $x$ e $y$ tiene la forma:
    \[
      \left\{
      \begin{array}{rcrcr}
        a_1 x & + & b_1 y & = & c_1 \\[6pt]
        a_2 x & + & b_2 y & = & c_2
      \end{array}
      \right.
    \]
  \end{block}

  \vspace{0.3em}
  \begin{columns}[T]
    \column{0.48\textwidth}
      \textbf{En contexto agronómico:}
      \begin{itemize}
        \item $x, y$: cantidades a determinar\\
          (dosis, volúmenes, superficies\ldots)
        \item $a_i, b_i$: coeficientes técnicos\\
          (concentraciones, rendimientos\ldots)
        \item $c_i$: requisitos del cultivo\\
          (demanda nutricional, restricciones\ldots)
      \end{itemize}
    \column{0.48\textwidth}
      \textbf{Interpretación geométrica:}
      \begin{itemize}
        \item Cada ecuación es una \textbf{restricción lineal}
          representada como una recta
        \item La solución es la \textbf{combinación factible}
          que cumple todas las restricciones simultáneamente
      \end{itemize}
  \end{columns}
\end{frame}

% ---- Forma matricial ---------------------------------------
\begin{frame}{Forma matricial}
  El sistema se escribe como $A\mathbf{x} = \mathbf{b}$:
  \[
    \underbrace{\begin{pmatrix} a_1 & b_1 \\ a_2 & b_2 \end{pmatrix}}_{A\;(\text{coeficientes técnicos})}
    \underbrace{\begin{pmatrix} x \\ y \end{pmatrix}}_{\mathbf{x}\;(\text{incógnitas})}
    =
    \underbrace{\begin{pmatrix} c_1 \\ c_2 \end{pmatrix}}_{\mathbf{b}\;(\text{requerimientos})}
  \]

  \vspace{0.3em}
  \begin{block}{Matriz aumentada — herramienta de cálculo}
    \[
      [A \mid \mathbf{b}] =
      \left(\begin{array}{cc|c}
        a_1 & b_1 & c_1 \\
        a_2 & b_2 & c_2
      \end{array}\right)
    \]
  \end{block}

  \vspace{0.2em}
  \begin{exampleblock}{Ejemplo nutricional}
    \[
      \left(\begin{array}{cc|c}
        0.20 & 0.05 & 30 \\
        0.10 & 0.15 & 20
      \end{array}\right)
      \leftarrow
      \begin{array}{l}
        \text{requerimiento de N (kg/ha)}\\
        \text{requerimiento de P (kg/ha)}
      \end{array}
    \]
  \end{exampleblock}
\end{frame}

% ===========================================================
\section{Clasificación de sistemas}
% ===========================================================

\begin{frame}{Clasificación: el determinante}
  Sea $D = \det(A) = a_1 b_2 - a_2 b_1$.

  \vspace{0.5em}
  \begin{columns}[T]
    \column{0.32\textwidth}
      \begin{exampleblock}{$D \neq 0$}
        \centering\textbf{Compatible\\determinado}\\[4pt]
        \textbf{Solución única}\\[4pt]
        Existe exactamente\\una combinación de\\insumos que cumple\\los requerimientos
      \end{exampleblock}

    \column{0.32\textwidth}
      \begin{block}{$D = 0$, consistente}
        \centering\textbf{Compatible\\indeterminado}\\[4pt]
        \textbf{Infinitas soluciones}\\[4pt]
        Las restricciones\\son redundantes:\\múltiples mezclas\\son equivalentes
      \end{block}

    \column{0.32\textwidth}
      \begin{alertblock}{$D = 0$, inconsistente}
        \centering\textbf{Incompatible}\\[4pt]
        \textbf{Sin solución}\\[4pt]
        Los insumos\\disponibles no pueden\\satisfacer todos los\\requerimientos
      \end{alertblock}
  \end{columns}
\end{frame}

% ---- Visualización geométrica ------------------------------
\begin{frame}{Visualización geométrica — restricciones de cultivo}
  \begin{columns}[c]
    \column{0.33\textwidth}
    \centering\small\textbf{Mezcla única posible}
    \begin{tikzpicture}[scale=0.75]
      \begin{axis}[
        xmin=0, xmax=5, ymin=0, ymax=5,
        axis lines=left,
        xlabel={\small $x$ (kg fert.\ A)},
        ylabel={\small $y$ (kg fert.\ B)},
        tick style={draw=none}, ticklabel style={font=\tiny},
        width=4.8cm, height=4.8cm,
        xlabel style={font=\tiny}, ylabel style={font=\tiny}
      ]
        \addplot[celesteagua, thick, domain=0:5] {(3-0.2*x)/0.05};
        \addplot[amarillotrigo, thick, domain=0:5] {(2-0.1*x)/0.15};
        \addplot[only marks, mark=*, mark size=2.5pt, verdecampo]
          coordinates {(100/14, 100/14)};
      \end{axis}
    \end{tikzpicture}

    \column{0.33\textwidth}
    \centering\small\textbf{Restricciones redundantes}
    \begin{tikzpicture}[scale=0.75]
      \begin{axis}[
        xmin=0, xmax=5, ymin=0, ymax=5,
        axis lines=left,
        xlabel={\small $x$}, ylabel={\small $y$},
        tick style={draw=none}, ticklabel style={font=\tiny},
        width=4.8cm, height=4.8cm,
        xlabel style={font=\tiny}, ylabel style={font=\tiny}
      ]
        \addplot[celesteagua, very thick, domain=0:5] {4-x};
        \addplot[amarillotrigo, thick, dashed, domain=0:5] {4-x};
      \end{axis}
    \end{tikzpicture}

    \column{0.33\textwidth}
    \centering\small\textbf{Requerimientos imposibles}
    \begin{tikzpicture}[scale=0.75]
      \begin{axis}[
        xmin=0, xmax=5, ymin=0, ymax=5,
        axis lines=left,
        xlabel={\small $x$}, ylabel={\small $y$},
        tick style={draw=none}, ticklabel style={font=\tiny},
        width=4.8cm, height=4.8cm,
        xlabel style={font=\tiny}, ylabel style={font=\tiny}
      ]
        \addplot[celesteagua, thick, domain=0:5] {3-x};
        \addplot[amarillotrigo, thick, domain=0:5] {5-x};
      \end{axis}
    \end{tikzpicture}
  \end{columns}

  \vspace{0.3em}
  \centering{\small\color{verdecampo}
    Azul: restricción de N \quad | \quad Amarillo: restricción de P}
\end{frame}

% ===========================================================
\section{Métodos de resolución}
% ===========================================================

\begin{frame}{Métodos de resolución}
  \begin{block}{Cuatro métodos equivalentes}
    \begin{enumerate}
      \item \textbf{Sustitución} — despejar una variable y reemplazar
      \item \textbf{Eliminación gaussiana} — operar ecuaciones para cancelar una incógnita
      \item \textbf{Regla de Cramer} — fórmula explícita mediante determinantes
      \item \textbf{Método matricial} — inversa de $A$ o Gauss-Jordan
    \end{enumerate}
  \end{block}

  \vspace{0.5em}
  \begin{exampleblock}{Sistema de referencia — Fertilización nitrogenada y fosfatada}
    Un técnico agrícola necesita aportar \textbf{30 kg/ha de N} y \textbf{20 kg/ha de P}
    combinando:
    \[
      \left\{
      \begin{array}{lcl}
        0.20\,x + 0.05\,y & = & 30 \quad \text{(N, kg/ha)} \\[4pt]
        0.10\,x + 0.15\,y & = & 20 \quad \text{(P, kg/ha)}
      \end{array}
      \right.
    \]
    donde $x$ = kg/ha de fertilizante A, \; $y$ = kg/ha de fertilizante B.
  \end{exampleblock}
\end{frame}

% ---- Sustitución -------------------------------------------
\subsection{Sustitución}

\begin{frame}{Método de sustitución — Fertilización N-P}
  \[
    \left\{
    \begin{array}{rcrcr}
      0.20\,x & + & 0.05\,y & = & 30 \\[4pt]
      0.10\,x & + & 0.15\,y & = & 20
    \end{array}
    \right.
  \]

  \begin{enumerate}
    \item<1-> Despejar $x$ de la primera ecuación:
      \[ x = \frac{30 - 0.05\,y}{0.20} = 150 - 0.25\,y \]

    \item<2-> Sustituir en la segunda:
      \[ 0.10(150 - 0.25\,y) + 0.15\,y = 20
         \;\Rightarrow\; 15 - 0.025\,y + 0.15\,y = 20 \]
      \[ 0.125\,y = 5 \;\Rightarrow\; \alert{y = 40 \unidad{kg/ha}} \]

    \item<3-> Retrosustituir:
      \[ x = 150 - 0.25 \times 40 = \alert{140 \unidad{kg/ha}} \]

    \item<4-> \textbf{Verificar}: $0.20(140) + 0.05(40) = 28 + 2 = 30$ ✓ \quad
      $0.10(140) + 0.15(40) = 14 + 6 = 20$ ✓

    \item<5-> \textbf{Interpretación}: aplicar \alert{140 kg/ha} del fertilizante A
      y \alert{40 kg/ha} del fertilizante B.
  \end{enumerate}
\end{frame}

% ---- Eliminación -------------------------------------------
\subsection{Eliminación gaussiana}

\begin{frame}{Método de eliminación — Balance hídrico}
  \begin{exampleblock}{Contexto}
    Un sistema de riego combina agua de pozo ($x$, m$^3$/h)
    y agua superficial ($y$, m$^3$/h). La conductividad eléctrica (CE)
    y la demanda total deben cumplir:
    \[
      \left\{
      \begin{array}{rcrcr}
        x  & + &  y  & = & 80 \quad \text{(caudal total, m}^3\text{/h)} \\[4pt]
        0.8\,x & + & 1.6\,y & = & 96 \quad \text{(CE ponderada, dS/m $\times$ m}^3\text{/h)}
      \end{array}
      \right.
    \]
  \end{exampleblock}

  \begin{enumerate}
    \item<2-> Multiplicar la primera ecuación por $0.8$:
      \[ 0.8\,x + 0.8\,y = 64 \]

    \item<3-> Restar de la segunda ecuación:
      \[ (0.8\,x + 1.6\,y) - (0.8\,x + 0.8\,y) = 96 - 64 \]
      \[ 0.8\,y = 32 \;\Rightarrow\; \alert{y = 40 \unidad{m}^3\text{/h}} \]

    \item<4-> Sustituir: $x = 80 - 40 = \alert{40 \unidad{m}^3\text{/h}}$

    \item<5-> \textbf{Interpretación}: mezclar a partes iguales para obtener
      CE $= 96/80 = 1.2$ dS/m.
  \end{enumerate}
\end{frame}

% ---- Cramer ------------------------------------------------
\subsection{Regla de Cramer}

\begin{frame}{Regla de Cramer}
  \begin{block}{Teorema}
    Si $D = \det(A) \neq 0$, la solución única del sistema es:
    \[
      x = \frac{D_x}{D}, \qquad y = \frac{D_y}{D}
    \]
    \[
      D = \deter{a_1}{b_1}{a_2}{b_2}, \quad
      D_x = \deter{c_1}{b_1}{c_2}{b_2}, \quad
      D_y = \deter{a_1}{c_1}{a_2}{c_2}
    \]
  \end{block}

  \pause
  \textbf{Aplicación — problema de fertilización:}
  \begin{align*}
    D   &= \deter{0.20}{0.05}{0.10}{0.15}
         = (0.20)(0.15)-(0.05)(0.10) = 0.030 - 0.005 = 0.025 \\[4pt]
    D_x &= \deter{30}{0.05}{20}{0.15}
         = (30)(0.15)-(0.05)(20) = 4.5 - 1.0 = 3.5 \\[4pt]
    D_y &= \deter{0.20}{30}{0.10}{20}
         = (0.20)(20)-(30)(0.10) = 4.0 - 3.0 = 1.0
  \end{align*}
  \pause
  \[
    x = \frac{3.5}{0.025} = \alert{140 \unidad{kg/ha}}, \qquad
    y = \frac{1.0}{0.025} = \alert{40 \unidad{kg/ha}}
  \]
\end{frame}

% ---- Matricial ---------------------------------------------
\subsection{Método matricial}

\begin{frame}{Método matricial — eliminación de Gauss-Jordan}
  Escribir la \textbf{matriz aumentada} y reducir a forma escalonada reducida:
  \[
    \left(\begin{array}{cc|c}
      0.20 & 0.05 & 30 \\
      0.10 & 0.15 & 20
    \end{array}\right)
    \xrightarrow{F_1 \leftarrow 5F_1;\; F_2 \leftarrow 10F_2}
    \left(\begin{array}{cc|c}
      1 & 0.25 & 150 \\
      1 & 1.5  & 200
    \end{array}\right)
  \]

  \pause
  \[
    \xrightarrow{F_2 \leftarrow F_2 - F_1}
    \left(\begin{array}{cc|c}
      1 & 0.25 & 150 \\
      0 & 1.25 & 50
    \end{array}\right)
    \xrightarrow{F_2 \leftarrow F_2 / 1.25}
    \left(\begin{array}{cc|c}
      1 & 0.25 & 150 \\
      0 & 1    & 40
    \end{array}\right)
  \]

  \pause
  \[
    \xrightarrow{F_1 \leftarrow F_1 - 0.25\,F_2}
    \left(\begin{array}{cc|c}
      1 & 0 & 140 \\
      0 & 1 & 40
    \end{array}\right)
    \Rightarrow
    \alert{x = 140 \unidad{kg/ha}, \quad y = 40 \unidad{kg/ha}}
  \]
\end{frame}

% ===========================================================
\section{Ejemplos resueltos}
% ===========================================================

\begin{frame}{Ejemplo 1 — Mezcla de plaguicidas}
  \begin{exampleblock}{Problema}
    Un técnico necesita preparar \textbf{200 L} de solución fungicida
    al \textbf{0.8\%} de ingrediente activo (i.a.).
    Dispone de dos concentrados:
    \begin{itemize}
      \item Producto X: 2\% i.a. \quad
      \item Producto Y: 0.5\% i.a.
    \end{itemize}
    ¿Cuántos litros de cada uno se deben mezclar?
  \end{exampleblock}

  \pause
  \textbf{Variables:} $x$ = L del producto X, \; $y$ = L del producto Y.

  \pause
  \textbf{Sistema:}
  \[
    \left\{
    \begin{array}{rcrcrl}
      x & + & y & = & 200 & \quad\text{(volumen total)}\\[4pt]
      0.02\,x & + & 0.005\,y & = & 0.008 \times 200 = 1.6 & \quad\text{(masa i.a.)}
    \end{array}
    \right.
  \]

  \pause
  \textbf{Solución} (eliminación: $F_2 \leftarrow F_2 - 0.005\,F_1$):
  \[
    0.015\,x = 1.6 - 1.0 = 0.6
    \;\Rightarrow\; \alert{x = 40\,\text{L}}, \quad \alert{y = 160\,\text{L}}
  \]
\end{frame}

\begin{frame}{Ejemplo 2 — Asignación de superficie entre cultivos}
  \begin{exampleblock}{Problema}
    Un agricultor dispone de \textbf{50 ha} para sembrar trigo y maíz.
    Su presupuesto es de \textbf{\$75\,000}.
    El costo de producción es \$1\,200/ha para trigo y \$1\,800/ha para maíz.
    ¿Cuántas hectáreas dedicar a cada cultivo?
  \end{exampleblock}

  \pause
  \textbf{Variables:} $t$ = ha de trigo, \; $m$ = ha de maíz.

  \[
    \left\{
    \begin{array}{rcrcr}
      t & + & m & = & 50 \\[4pt]
      1200\,t & + & 1800\,m & = & 75000
    \end{array}
    \right.
  \]

  \pause
  \textbf{Cramer:}
  \[
    D = \deter{1}{1}{1200}{1800} = 600, \quad
    D_t = \deter{50}{1}{75000}{1800} = -\!15000, \quad
    D_m = \deter{1}{50}{1200}{75000} = 15000
  \]
  \pause
  \[
    t = \frac{-15000}{600}= \alert{25\,\text{ha de trigo}}, \qquad
    m = \frac{15000}{600} = \alert{25\,\text{ha de maíz}}
  \]
\end{frame}

\begin{frame}{Ejemplo 3 — Sistema incompatible en nutrición vegetal}
  \begin{exampleblock}{Contexto}
    Se pretende cubrir \textbf{40 kg/ha de N} y \textbf{20 kg/ha de P} con dos
    fertilizantes que tienen la misma relación N:P = 2:1.
    \[
      \left\{
      \begin{array}{rcrcr}
        0.20\,x & + & 0.10\,y & = & 40\\[4pt]
        0.10\,x & + & 0.05\,y & = & 20
      \end{array}
      \right.
    \]
  \end{exampleblock}

  \begin{enumerate}
    \item<2-> $D = \deter{0.20}{0.10}{0.10}{0.05} = 0.010 - 0.010 = 0$
    \item<3-> La segunda ecuación es exactamente la mitad de la primera:
      son \textbf{linealmente dependientes}.
    \item<4-> Las dos ecuaciones representan la \textbf{misma restricción},
      por lo que el sistema es \textbf{compatible indeterminado}:
      existe infinidad de combinaciones $x, y \geq 0$ que satisfacen ambas.
    \item<5->
      \begin{alertblock}{Interpretación agronómica}
        Si dos fertilizantes tienen la misma proporción de nutrientes,
        no se puede fijar simultáneamente la dosis de N \textit{y} la de P
        con una sola restricción adicional. Se necesita un tercer criterio
        (p.ej., costo o disponibilidad).
      \end{alertblock}
  \end{enumerate}
\end{frame}

% ===========================================================
\section{Aplicaciones adicionales}
% ===========================================================

\begin{frame}{Aplicación — Rotación de cultivos y rentabilidad}
  \begin{exampleblock}{Problema}
    Una explotación de \textbf{120 ha} puede destinar su superficie a papa o cebolla.
    La ganancia esperada es \$3\,500/ha en papa y \$2\,000/ha en cebolla.
    Para optimizar el uso de mano de obra, la ganancia total debe ser \$\mathbf{330\,000}$.
    ¿Cuántas hectáreas asignar a cada cultivo?
  \end{exampleblock}

  \pause
  \[
    \left\{
    \begin{array}{rcrcrl}
      p & + & c & = & 120 & \quad \text{(ha totales)} \\[4pt]
      3500\,p & + & 2000\,c & = & 330000 & \quad \text{(ganancia objetivo)}
    \end{array}
    \right.
  \]

  \pause
  Sustitución: $p = 120 - c$ en la segunda:
  \[
    3500(120-c) + 2000\,c = 330000
    \;\Rightarrow\; 420000 - 1500\,c = 330000
    \;\Rightarrow\; \alert{c = 60\,\text{ha}}
  \]
  \pause
  \[
    p = 120 - 60 = \alert{60\,\text{ha}}
  \]

  \vspace{0.2em}
  \textbf{Verificación}: $3500(60) + 2000(60) = 210000 + 120000 = 330000$ ✓
\end{frame}

\begin{frame}{Aplicación — Preparación de sustratos}
  \begin{exampleblock}{Problema}
    Un viverista prepara \textbf{1\,000 L} de sustrato mezclando:
    \begin{itemize}
      \item Turba: densidad aparente $0.12\,\text{g/cm}^3$, pH 3.8
      \item Perlita: densidad aparente $0.08\,\text{g/cm}^3$, pH 7.0
    \end{itemize}
    El sustrato final debe pesar \textbf{100 kg} y tener pH \textbf{5.0} (ponderado por volumen).
  \end{exampleblock}

  \pause
  \textbf{Variables:} $t$ = L de turba, \; $p$ = L de perlita.
  \[
    \left\{
    \begin{array}{rcrcrl}
      t & + & p & = & 1000 & \quad\text{(volumen, L)} \\[4pt]
      0.12\,t & + & 0.08\,p & = & 100 & \quad\text{(masa, kg)}
    \end{array}
    \right.
  \]

  \pause
  Eliminación ($F_2 \leftarrow F_2 - 0.08\,F_1$):
  \[
    0.04\,t = 100 - 80 = 20
    \;\Rightarrow\; \alert{t = 500\,\text{L}}, \quad \alert{p = 500\,\text{L}}
  \]

  \pause
  Verificación del pH: $0.5 \times 3.8 + 0.5 \times 7.0 = 1.9 + 3.5 = 5.4$ \\
  \begin{alertblock}{Nota práctica}
    El pH ponderado resultante (5.4) difiere del objetivo (5.0).
    La restricción de pH requeriría una tercera variable (enmienda).
  \end{alertblock}
\end{frame}

% ===========================================================
\section{Resumen y conclusiones}
% ===========================================================

\begin{frame}{Cuadro comparativo de métodos}
  \begin{table}
    \centering
    \renewcommand{\arraystretch}{1.4}
    \begin{tabular}{llll}
      \toprule
      \textbf{Método} & \textbf{Ventaja} & \textbf{Limitación}
        & \textbf{Uso típico} \\
      \midrule
      Sustitución  & Intuitivo       & Tedioso con decimales
        & Mezclas simples \\
      Eliminación  & Sistemático     & Requiere cuidado con fracciones
        & Balance hídrico \\
      Cramer        & Fórmula directa & Solo si $D \neq 0$
        & Fertilización \\
      Matricial    & Generalizable   & Requiere álgebra matricial
        & Software técnico \\
      \bottomrule
    \end{tabular}
  \end{table}
\end{frame}

\begin{frame}{Resumen}
  \begin{block}{Ideas clave}
    \begin{enumerate}
      \item Muchos problemas agronómicos implican \textbf{dos condiciones simultáneas}
        que se modelan con sistemas $2\times2$.
      \item El determinante $D$ indica si el problema tiene solución única,
        infinitas o ninguna --- lo cual tiene \textbf{significado físico real}.
      \item Todos los métodos son algebraicamente equivalentes;
        se elige según la \textbf{complejidad numérica} del problema.
      \item La \textbf{verificación} es indispensable: en agronomía,
        un error de cálculo puede arruinar una cosecha o generar un impacto ambiental negativo.
      \item La forma matricial $A\mathbf{x}=\mathbf{b}$ permite escalar a
        \textbf{$n$ nutrientes y $n$ fuentes} (sistemas $n\times n$).
    \end{enumerate}
  \end{block}

  \vspace{0.3em}
  \begin{exampleblock}{Próxima clase}
    Sistemas $3\times3$ aplicados a nutrición vegetal con tres macronutrientes
    (N, P, K) y tres fuentes de fertilizante.
  \end{exampleblock}
\end{frame}

% ---- Ejercicios propuestos ----------------------------------
\begin{frame}{Ejercicios propuestos}
  \small
  \begin{columns}[T]
    \column{0.50\textwidth}
    \begin{enumerate}
      \item \textbf{Fertirriego}: una solución nutritiva requiere
        $8\,\text{meq/L}$ de Ca$^{2+}$ y $4\,\text{meq/L}$ de Mg$^{2+}$.
        Usando nitrato de calcio (15-0-0 + 19\% Ca) y sulfato de magnesio
        (0\% N + 10\% Mg), plantear y resolver el sistema.

      \vspace{0.6em}
      \item \textbf{Siembra mixta}: un campo de 80 ha se siembra con soja
        y girasol. El agua disponible es 12\,000 m$^3$ y la demanda
        hídrica es 180 m$^3$/ha para soja y 120 m$^3$/ha para girasol.
        ¿Cuántas hectáreas de cada cultivo?
    \end{enumerate}

    \column{0.48\textwidth}
    \begin{enumerate}
      \setcounter{enumi}{2}
      \item \textbf{Compostaje}: se mezclan dos residuos para obtener
        500 kg de compost con relación C:N = 25.
        Residuo A: C:N = 30, Residuo B: C:N = 20.
        Plantear el sistema y clasificarlo.

      \vspace{0.6em}
      \item \textbf{Clasificación}: determinar si el siguiente sistema
        de restricciones hídricas tiene solución y justificar
        agronómicamente:
        \[
          \left\{\begin{array}{l}
            2x + 4y = 800 \\
            x + 2y = 500
          \end{array}\right.
        \]
    \end{enumerate}
  \end{columns}
\end{frame}

% ---- Fin ---------------------------------------------------
\begin{frame}[plain]
  \begin{center}
    {\Large\color{verdecampo}\textbf{¡Gracias!}}\\[1em]
    {\normalsize Preguntas y discusión}\\[2em]
    \textit{``La agricultura es la profesión más antigua y más noble del mundo;
    las matemáticas, su lenguaje oculto.''}\\[1.5em]
    \begin{tikzpicture}
      \fill[verdecampo!30, rounded corners=4pt] (0,0) rectangle (9,0.25);
      \fill[verdeagro!50, rounded corners=4pt] (0,0) rectangle (6,0.25);
      \fill[amarillotrigo!80, rounded corners=4pt] (0,0) rectangle (3,0.25);
    \end{tikzpicture}
  \end{center}
\end{frame}

\end{document}
